\documentclass[a4paper,12pt,twoside]{article}
\usepackage[utf8]{inputenc}
\usepackage[T1]{fontenc}
\usepackage[finnish]{babel}
\usepackage{lmodern}
\usepackage{eurosym}
\usepackage[pdftex,colorlinks]{hyperref}
\begin{document}
\title{Tiedostopohja}
\author{Essi Esimerkki}
\maketitle
\tableofcontents
\newtheorem{esimerkki}{Esimerkki}
\section{Lorem ipsum}
Lorem ipsum dolor sit amet, consectetur adipiscing elit. Pellentesque ac purus sit amet mauris dapibus facilisis sed a eros. Vestibulum vel interdum eros. Aenean vehicula eget nisl sed porttitor. Donec urna mi, aliquet ac malesuada quis, ultricies at purus. Curabitur a ligula nisi. Morbi arcu leo, cursus a orci nec, facilisis egestas urna. Sed tempor, nulla et cursus sagittis, mauris nibh dignissim purus, quis iaculis lectus turpis quis libero. Vivamus ut tortor non erat gravida egestas \cite{Aho}. Cras ut ligula quis felis consectetur bibendum. Nulla aliquet placerat turpis, nec rutrum metus.\footnote{Tämä on laskelmoitua.} Lorem ipsum dolor sit amet, consectetur adipiscing elit.
\subsection{Kaikki tosiseikat}
Proin fringilla tristique ipsum a feugiat. Curabitur eros enim, vulputate id mauris vel, tincidunt viverra tortor (\ref{vaihtoehtoja}). Vivamus convallis, felis at lobortis volutpat, leo urna ultrices elit, et volutpat nisi arcu non diam. Vivamus sed eleifend lectus, quis ultricies felis. Nam hendrerit ornare tellus ut vehicula [Esimerkki \ref{meow}]. Praesent ullamcorper tincidunt facilisis. Pellentesque facilisis quis ligula vel rhoncus. Donec posuere tellus ac urna maximus blandit. Lorem ipsum dolor sit amet, consectetur adipiscing elit \cite{Ber}. Integer commodo viverra felis sit amet blandit. Sed ut tincidunt lectus. Vivamus sed scelerisque urna. Sed egestas molestie justo id mattis.
\begin{equation}
\label{vaihtoehtoja}
f(n+1,\vec{x})=
\left\{
\begin{array}{cl}
g(\vec{x}), & n=0 \\
h\left(n,f(n,\vec{x}),\vec{x}\right) & \textrm{muuten}
\end{array}
\right.
\end{equation}
\subsection{Lyhyt}
100 \euro
\begin{esimerkki}[Meow]
\label{meow}
Kissojen Javascript-koodausta.
\end{esimerkki}
\begin{thebibliography}{99}
\bibitem[Aho]{Aho} Aho: Koodausopas
\bibitem[Ber]{Ber} Berg: Kissakuvasto
\end{thebibliography}
\end{document}
